\[
  3a+1-2e=2^jJ(y),\qquad j\ge2,
\]
so $y=B(x)$. Define
\[
 P=Q_1^e(a),\quad U=Q_2^e(a),\quad
 D=Q_j^g(J(y)),\quad C=Q_{j-1}^g(J(y)),\quad F=A(U).
\]
Then $D=A(P)$ and $C=B(P)=B(U)$. If $X$ is the common child of $D,C$ and
$Y$ is the other child of $C$, then
\[
  B(F)=Y.
\]
Consequently an obligation in the $B$-selecting phase transfers a continuing
\DRAW\ to the adjacent frame $\{D,C\}$; when $j=2$ the transferred frame is
again an obligation in the opposite phase.
\end{lemma}

\begin{proof}
The formulas for $D$ and $C$ are the signed boundary formula and the
shared-child identity. The remaining equality is a direct substitution in the
three cases $j=2$, $j=3$, and $j\ge4$: the constant suffix of $D$ determines
its common child $X$, while the one-step longer factor state $F$ has the same
raw side $Y$. For the outcome assertion, a common child of a \DRAW\ cannot be
\LOSS. If the exponent-one member is \DRAW, the frame $\{D,C\}$ already
contains a \DRAW. In the parent form of an obligation, the only remaining
orientation has the exponent-two member \DRAW; if $C$ were \WIN, outcome
recursion through the displayed diamond would force $Y$ both \LOSS\ and a
child of the \DRAW\ state $F$, impossible. Hence $C$ is \DRAW. For $j=2$,
$D,C$ are exactly the exponent-two/one pair over $J(y)$.
\end{proof}


\section{The closed routing bridge}\label{sec:closed-routing}

The local arithmetic in Sections 2--8 is deliberately left unchanged from the
audit-hardened reduction.  We now close the two global obligations that were
left open there.  Two negative regressions are retained throughout:

\begin{enumerate}[label=(N\arabic*)]
\item a smaller canonical coefficient is not, in the presence of a factor
      $3^k$, a smaller coefficient source;
\item a $B$ return after several canonical $A$ steps need not lie below the
      numerical source present before that streak.
\end{enumerate}

The proof below never uses either implication.  Numerical values displayed
along a canonical streak are routing cursors until a stated source comparison
promotes one of them to a ranked anchor.

\subsection{Permanent multi-$A$ regression}

\begin{lemma}[Pre-streak anchor counterexample]\label{lem:anchor-stress}
Let $x_0=20$ and $x_{i+1}=A(x_i)$.  Then
\[
 20\longmapsto30\longmapsto45\longmapsto68,
 \qquad B(68)=25>20.
\]
The phase carried by the canonical lift switches to the $B$-selecting phase
at $68$, and the corresponding signed valuation is $3$.
\end{lemma}

\begin{proof}
The three $A$ values are immediate.  Moreover
\[
 \alpha(20),\alpha(30),\alpha(45),\alpha(68)=1,0,1,1,
\]
whereas the lifted phase flips after each canonical $A$ continuation.  At
$68$ the carried phase is therefore $0\ne\alpha(68)$, and
\[
 3J(68)+1=616=2^3\,77,
 \qquad 77=J(25).
\]
Thus the return is $25$, not a value below the pre-streak source $20$.
\end{proof}

\subsection{The exact raw selector}

For $s>0$ and $e\in\{0,1\}$ define the ordinary source selected by phase $e$
by
\[
 \chi_e(s)=
 \begin{cases}
   A(s),&e=\alpha(s),\\
   B(s),&e=1-\alpha(s).
 \end{cases}
\]

\begin{lemma}[Raw-selector identity]\label{lem:raw-selector}
For every $s>0$ and phase $e$,
\[
 \boxed{R\!\left(Q_1^e(J(s))\right)=\chi_e(s).}
\]
\end{lemma}

\begin{proof}
Put $z=Q_1^e(J(s))=2J(s)-e$ and $t=R(z)$.  Since
$z\bmod2=e$, Lemma~\ref{lem:alt-lift}, applied to $z$, gives for some
$m\ge1$
\[
 3z+1=Q_m^{1-e}(J(t))=2^mJ(t)-(1-e).
\]
After substituting $z=2J(s)-e$ and dividing by $2$,
\[
 3J(s)+1-2e=2^{m-1}J(t).
\]
The unique factor-free source in this signed transition is therefore $t$.
Lemma~\ref{lem:source-boundary} says that this source is $A(s)$ in the
$A$-selecting phase and $B(s)$ in the opposite phase.  Hence
$t=\chi_e(s)$.
\end{proof}

The next lemma is the missing provenance statement behind the old
$A/B_2$ reset.

\begin{lemma}[Second selector across an ordinary side diamond]
\label{lem:second-selector}
Let $q>0$ be ordinary, put $y=B(q)$ and $g=\alpha(q)$.  Then
\[
 \boxed{B(A(q))=\chi_g(y).}
\]
In particular, if the phase at $q$ is $B$-selecting and its signed valuation
is $2$, the transferred obligation at $y=B(q)$ selects on its next source
step the exact common side $B(A(q))$.
\end{lemma}

\begin{proof}
Let $k$ be the alternating-suffix length of $A(q)$.  Ordinary means
$k\in\{1,2\}$.  The proof of Lemma~\ref{lem:side} gives the refined table
\[
\begin{array}{c|c|c|c}
 k&g& B(A(q))=A(y) & g=\alpha(y)\\ \hline
 1&1-(y\bmod2)&y\bmod4\in\{0,3\}&y\bmod4\in\{0,3\}\\
 2&y\bmod2&y\bmod4\in\{1,2\}&y\bmod4\in\{1,2\}.
\end{array}
\]
The formulas for $g$ follow directly by substituting the one-bit and two-bit
alternating suffixes in $A(q)=2^k y+t_{k,y\bmod2}$ and using the two-bit
definition of $\alpha$.  Thus $B(A(q))=A(y)$ exactly when
$g=\alpha(y)$; otherwise it is $B(y)$.  Lemma~\ref{lem:raw-selector}
identifies this with $\chi_g(y)$.

For the last assertion note that the signed $B$-valuation equals $k+1$.
Valuation $2$ therefore has $k=1$, so $q$ is ordinary automatically; the
transferred phase is $g=1-(1-\alpha(q))=\alpha(q)$.
\end{proof}

\begin{corollary}[Exceptional rows never form a $B_2$ recycle]
\label{cor:exceptional-no-b2}
If $q\bmod16\in\{1,3,12,14\}$, then the $B$-selecting signed valuation at
$q$ is at least $4$.
\end{corollary}

\begin{proof}
For an exceptional $q$, the alternating suffix of $A(q)$ has length at least
three.  The signed $B$-valuation is one more than that suffix length, hence is
at least four.
\end{proof}

\subsection{Retained data, two one-shot seeds, and phase alignment}

The earlier versions used two different finite-token layers.  That distinction
is essential and is restored here, but without the invalid pre-streak numerical
comparison of Lemma~\ref{lem:anchor-stress}.  A marked configuration retains

\[
  (s,\eta,M,\zeta,N,a;\ \text{routing data}).
\]

Here $s$ is the outer coefficient-source anchor; $M$ is a finite multiset of
certified outer proof-token heights; $\eta\in\{1,0\}$ is the unique outer
seed bit; $N$ is a finite multiset of certified inner proof-token heights;
$\zeta\in\{1,0\}$ is the inner seed bit; and $a$ is the retained numerical
anchor supplied by the currently active typed entry.  When the entry is a
free obligation its factor-free coefficient is $J(a)$.  In an attached
factor/lift task, $a$ is the nested source anchor inherited from the exact
parent/child geometry.  Temporary ordinary sources, phases, valuations and
tail exponents are routing data and are not silently promoted to $s$ or $a$.

Put
\[
 \Phi(M)=\mathop{\#}_{h\in M}\omega^h,
 \qquad
 \Psi(N)=\mathop{\#}_{h\in N}\omega^h,
\]
and order retained data lexicographically by
\begin{equation}\label{eq:v6-rank}
 \boxed{\Pi=(s,\eta,\Phi(M),\zeta,\Psi(N),a).}
\end{equation}
Both bits are ordered by $1>0$.  A step classified as an outer token
decrease preserves the preceding outer source $s$; a step classified as an
inner token decrease preserves $(s,\eta,M,\zeta)$.  Thus no decrease of
$\Phi(M)$ or $\Psi(N)$ can pay for an increase in an earlier retained
component.  The seed bits are governed by the following non-reseeding rule.

\begin{lemma}[Seed admissibility and non-reseeding]\label{lem:seed-policy}
At a new outer-source fibre, while $\eta=1$, the inner seed $\zeta$ is
unavailable.  The first exact finite WIN/LOSS token produced by any typed
entry diamond---boundary, loss-sibling, base-entry, or the first free
obligation---may be installed in $M$ only together with the unique strict
transition $\eta:1\to0$.  After that transition, $M$ changes only by
certified proper-descendant replacement.

At fixed $(s,\eta=0,M)$, the unique transition $\zeta:1\to0$ may install the
first exact finite token of the active attached inner task in $N$.  After
$\zeta=0$, $N$ changes only by certified proper-descendant replacement.  If
both seed bits have been spent, no unrelated finite state may ever be added
to either multiset merely because its outcome is known.  This restriction is
about \emph{ranked tokens}.  A typed fixed-fibre subroutine may use a finite
WIN/LOSS state of unrelated height as temporary routing data, provided it is
not inserted into $M$ or $N$ and the subroutine is proved well-founded for
every possible initial value of that temporary witness.  This is the well-founded-fibre distinction used later for attached fixed-fibre subroutines.

A strict decrease of $M$ may reset $\zeta,N,a$; a strict decrease of $N$ may
reset $a$.  With all token data fixed, $a$ changes only by a separately
proved strict source comparison.  In particular, the two seed transitions
cannot occur in the opposite order and neither can be repeated in an
unchanged earlier fibre.
\end{lemma}

\begin{proof}
This is exactly the lexicographic discipline of \eqref{eq:v6-rank}.  Since
$\eta$ precedes $\Phi(M)$, the transition $\eta:1\to0$ is strict even when it
installs an arbitrary finite first multiset $M$.  While $\eta=1$, allowing a
$\zeta$ transition would leave the earlier unused seed available and would
permit two incomparable initializations, so it is excluded by definition of
the marked relation.  Once $\eta=0$, $\zeta$ precedes $\Psi(N)$ and gives the
same one-shot payment for the first inner token.  After a seed is zero, adding
an unrelated token would increase the corresponding ordinal multiset and is
therefore not a permitted transition.  Only a strict decrease in an earlier
component permits later components to be reset.
\end{proof}

\begin{lemma}[Phase-aligned common side]\label{lem:phase-aligned-side}
Let $O(x,e)$ be an $A$-selecting obligation, so $e=\alpha(x)$.  Put
\[
 y=A(x),\qquad g=1-e,
\]
and let $b$ be the common side of its exponent-one/two boundary.  In the
canonical branch, whose next obligation is $O(y,g)$,
\[
 \boxed{b=\chi_g(y).}
\]
Consequently, if $g$ is $A$-selecting at $y$, then $b=A(y)$; if $g$ is
$B$-selecting, then $b=B(y)$.
\end{lemma}

\begin{proof}
With the notation of the obligation put
\[
 P=Q_1^e(J(x)),\qquad V=A(P).
\]
The canonical identity is
\[
 V=Q_1^g(J(y)).
\]
The common side is $b=B(P)=R(A(P))=R(V)$.  Lemma~\ref{lem:raw-selector}
therefore gives $b=\chi_g(y)$ exactly.
\end{proof}

\begin{lemma}[Finite-source front end]\label{lem:finite-front}
Suppose an $A$-selecting obligation $O(x,\alpha(x))$ is reached while the
ordinary source $x$ itself is already a retained finite WIN/LOSS token.
Before any later high return or free reset, one of the following occurs:
\begin{enumerate}[label=(\alph*)]
\item that token is replaced by a certified proper proof-tree descendant;
\item if $x$ is exceptional and its canonical LOSS child is $B(x)$, the
      route becomes an exact adjacent factor-free lift over this strict
      descendant token.
\end{enumerate}
The statement includes both possible outcomes, WIN or DRAW, of the common
side; no fresh finite endpoint of unrelated height is inserted.
\end{lemma}

\begin{proof}
Put $y=A(x)$ and let $b$ be the common side of the obligation.  As in the
universal obligation fork, $b$ is never LOSS.

If $x$ is LOSS, then $y$ is WIN and $h(y)\le h(x)-1$.  Since one child of
the WIN state $y$ is the non-LOSS state $b$, its other child $\ell$ is the
unique LOSS child.  Hence
\[
 h(\ell)=h(y)-1\le h(x)-2.
\]
Thus $x\mapsto\ell$ is strict, independently of whether $b$ is WIN or DRAW.

Now let $x$ be WIN and choose a canonical LOSS child $\ell_0$ of $x$.
If $\ell_0=y=A(x)$, then every child of $y$ is WIN; in particular the common
side $b$ is WIN and
\[
 h(b)\le h(\ell_0)-1=h(x)-2.
\]
So $x\mapsto b$ is strict.

It remains that $\ell_0=B(x)$.  If $x$ is ordinary, the ordinary side
identity makes
\[
 p=B(A(x))
\]
a child of the LOSS state $\ell_0$.  Therefore $p$ is WIN and
\[
 h(p)\le h(\ell_0)-1=h(x)-2.
\]
This strict token is retained before the obligation chooses any later
canonical or factor continuation.

If $x$ is exceptional, the unconditional exceptional child-pair identity
(Lemma~\ref{lem:exceptional-side}) gives the two children of $y=A(x)$ as an
exact adjacent pair
\[
 \{A(y),B(y)\}
 =\{Q_{m+1}^{\delta}(J(\ell_0)),Q_m^{\delta}(J(\ell_0))\}
\]
for some $m\ge1$ and phase $\delta$.  The token replacement
$x\mapsto\ell_0$ is strict because $h(\ell_0)=h(x)-1$, and the later
arithmetic is attached to this same token.  These cases are exhaustive.
\end{proof}

\begin{lemma}[Aligned obligation catches or spends its token]
\label{lem:aligned-catch}
Suppose $O(x,e)$ carries a retained WIN token
\[
 w=\chi_e(x).
\]
Then the next outcome-compatible macro either reaches an obligation whose
ordinary source is exactly the retained token $w$, or strictly replaces
$w$ by a proper descendant before any reset.

If $e$ is $B$-selecting, every valuation-two return reaches source $w$
exactly; valuations at least three enter an attached factor frame over $w$.
If $e$ is $A$-selecting, then $w=A(x)$.  A canonical continuation reaches
source $w$; every noncanonical continuation replaces $w$ by its unique LOSS
child before entering the attached factor/side constructor.
\end{lemma}

\begin{proof}
If $e$ is $B$-selecting, the selected ordinary source is by definition
$w=B(x)$.  Lemma~\ref{lem:b-diamond} transfers the DRAW to the exact adjacent
frame over $w$.  At valuation two this frame is the next obligation itself,
with ordinary source $w$; at larger valuation it is an attached factor frame
over the same retained token.

Now let $e=\alpha(x)$.  Then $w=A(x)$.  Let $d$ be the common side of
$O(x,e)$.  Lemma~\ref{lem:phase-aligned-side}, applied one level later,
places $d$ among the two ordinary children of $w$.  The state $d$ is never
LOSS.  If the obligation takes its canonical branch, the next ordinary
source is $A(x)=w$ exactly.  Otherwise $d$ is either DRAW or the WIN common
side of the factor branch.  Since the retained state $w$ is WIN and one of
its children $d$ is non-LOSS, its other child $\ell$ is the unique LOSS
child, so
\[
 h(\ell)=h(w)-1.
\]
Replace $w$ by $\ell$ before entering the corresponding attached side/factor
constructor.  Thus no noncanonical row can reset at equal token rank.
\end{proof}

\subsection{Closure of the seeded $A/B_2$ reset (former P1)}

\begin{proposition}[Seeded $A/B_2$ reset provenance]\label{prop:P1-closed}
Every outcome-compatible canonical $A$ continuation and every valuation-two
$B$ return from a typed free/aligned obligation satisfies the P1 acceptance
conditions.  More precisely, before a control reset to a fresh $A$-obligation
is permitted, one of the following has happened:
\begin{enumerate}[label=(\roman*)]
\item a retained numerical source has strictly decreased;
\item an available seed has installed an exact finite token according to
      Lemma~\ref{lem:seed-policy};
\item an existing retained proof token has been replaced by a certified proper
      proof-tree descendant; or
\item the route has become an attached factor/lift task over such a retained
      token, in which case it is no longer eligible for a free reseed.
\end{enumerate}
This statement covers canonical $A$-streaks of every outcome-compatible
length and all exceptional residue classes.  It never compares a final
$B$-return with the numerical cursor present before the streak.
\end{proposition}

\begin{proof}
First consider a free obligation whose ordinary source $x$ is the current
retained inner anchor and which does not yet carry an aligned finite token.
Such a free obligation is admissible only while a seed is available under
Lemma~\ref{lem:seed-policy}; once both seeds have been spent, all obligations
arrive through the aligned/attached cases below.  A $B$-selecting source
returns to $B(x)<x$, hence makes a genuine retained inner-source decrease.
At an $A$-selecting source the universal obligation fork has a common side
$b$ which is not LOSS.  If $b$ is DRAW, the standard source estimate
\[
 \rho(b)\le\frac{9x+1}{24}<x
\]
is a genuine source comparison against the factor-free retained source $x$
and gives (i).  If $b$ is WIN, install this exact boundary token using the
earliest admissible seed: $\eta:1\to0$ if this is the first finite token in
the outer fibre, otherwise $\zeta:1\to0$ with $\eta=0$ for the first inner
token of the current task.  No such installation is allowed if both bits are
zero.

Suppose first that the obligation then takes its factor branch.  The token
$b$ remains marked, and the factor fork is by definition an attached entry
over this exact token.  This is alternative (iv), so P1 stops here: it makes
